\begin{center}\bfseries Abstract\end{center}
\begingroup
\fontsize{9.5}{10.7}\selectfont
\setlength{\parskip}{1.5pt plus .5pt minus .5pt}
\noindent
A positive mass operator is constructed on the fibres of persistent
particles, and its scalar specialisation is determined for the central
electron. The result connects particle identity, state classification,
and mass evaluation: the magnitude arises from a structure preserving
the route, orientation, and memory of its generation. In the electronic
sector, two bidirectional readers agree on an explicit trajectory and
determine the same return correction; the remaining collections retain the
operator, spectral, or scalar scope proved for their own fibres.

The central trajectory visits exactly three nonadic diagonals and twenty-seven
cells. Two histories with the same absolute-value histogram retain distinct
reversible memories. The finite transform of the return word isolates three
nonzero modes and proves spectrally that $(\Omega,P_6)=(80,6)$. On the multisections, the
relative mass law yields a necessary and sufficient operator criterion: a
bijective TPK transport preserves mass exactly when it commutes with the mass
operator, equivalently when it preserves its complete logarithmic coordinate.

The regional census runs through \(104{,}976\) seed pairs, produces \(468\)
emissions, and retains \(243\) distinct words. Pairwise trit grouping places
those words in the \(729\)-position body \(S_0=E_9^3\). The exact threshold
difference then determines the shell \(243+27+1=271\) and the completion
\(729+243+27+1=1000\), with a normalisation compatible under the projections
between dimensions. Regional selection acts afterwards: neither the bijection
nor the cubic completion defines a species or a metrological value.

The proposed framework is Triadic Modular Holography, henceforth HMT, and
its realisation through Dimensional Mechanics. Its kernel comprises three
objects: \emph{All Possible Paths}, henceforth APP, is an arithmetic path
structure on \((\mathbb Z/9\mathbb Z)^2\): its additive Cayley graph,
with cardinal generators, fixes adjacency, and its cellular realisation
is toroidal. Sum and product operations and evaluations are distinguished
within that structure; TRIT is the unit
of oriented signature \(\{-1,0,+1\}\), distinguishing local regimes;
the \emph{Topological Phase Kernel}, henceforth TPK, composes the dynamics
of selection, transport, and memory. Its enriched evolution preserves
sheets, remainders, quotients, boundary, and carry, and produces the joint
discrete structure of the continuum. The composition of \(108\) updates
determines the dodecaphase record, and the orbital transform reversibly
reconstructs \(K\),
antecedents of the coordinates used by the mass reader.

A particle is defined through the persistence of compatible extensions;
the route records its observable history. The atlas distinguishes \(81\)
cells, \(56\) route configurations, \(13\) multisections, and \(324\) oriented
conditions. Their quotients, flavour and colour observables, the unique
central section, and the operations of conjugation, helicity, and holonomy
are constructed, including the Möbius band with its fibre and loop. The
integral signature and channel moments determine the positive character;
action, the clock, and the unit chart supply the absolute section. Mass
ratios preserve the relative character and differences between return
exponents.

The marked longitudinal incidence first fixes
\(N=5760\), \(r_N=5759/23040\), and
\(L_\alpha=\alpha_{\rm HMT}^{16}r_NL_\star\). That length induces the
clock, \(R_{108}=L_\alpha\), the energy \(Q_L\), the clock mass, and the
cotransformation of the electronic factor. On every massive fibre, the
identities \(H\Lambda=\hbar c_{\rm int}I\) and
\(\mathcal R_g\Lambda=L_\alpha^2I\) uniquely determine
\(G=c_{\rm int}^3L_\alpha^2/\hbar\); its circular expression is a
subsequent corollary. The deduction neither uses \(G\) as an antecedent
of the mass operator nor alters the spectrum tables.

Composition takes place on records before their signatures are evaluated.
The development encompasses zero-mass sectors, oriented propagation,
mixing, resonances, and Kepler--Sommerfeld orbital realisations. Graded
completion produces occupation rules and Pauli exclusion on complete
modes, with their spin, sheet, flavour, and colour labels. Metrological
comparison brings together each evaluation with its route, operator,
and chart, distinguishing central values, intervals, bounds, and widths.
Structure, equation, and value are thereby related while preserving the
specific content of each result and its physical identification.

\noindent\textbf{Keywords:} derivation of the mass spectrum; particle
structure; Triadic Modular Holography; Dimensional Mechanics; persistence;
holonomy; conjugation; metrological comparison.
\par\endgroup
